\documentclass[12pt]{article}

% Estrutura do documento.
\usepackage{geometry}
 \geometry{
 a4paper,
 total={170mm,257mm},
 left=25mm,
 top=25mm,
 right=20mm
 }
 
% Permite que o texto fique justificado
\usepackage[document]{ragged2e}

% Permite colocar palavras com acentos e configura como língua padrão o português
\usepackage[utf8]{inputenc}
\usepackage[brazilian]{babel}

% Permite carregar imagens no documento.
\usepackage{graphicx}
\usepackage{indentfirst}

% Fontes...
%\usepackage{bookman}
%\usepackage{boisik}
\usepackage{kpfonts}

% Para definir o espaçamento entre as linhas
\usepackage{setspace}

% Permite mostrar trechos de códigos/pseudocódigo no documento
\usepackage{listingsutf8}
\usepackage{color}

% Algumas cores
\definecolor{dkgreen}{rgb}{0,0.6,0}
\definecolor{lgray}{rgb}{0.960, 0.960, 0.960}
\definecolor{lblue}{rgb}{0.184, 0.552, 0.654}
\definecolor{gray}{rgb}{0.5,0.5,0.5}
\definecolor{mauve}{rgb}{0.58,0,0.82}
\definecolor{dkblue}{rgb}{0.109, 0.125, 0.592}

% Define como serão exibidos código no documento.
\lstset{frame=lines,
  inputencoding=utf8/latin1,
  aboveskip=3mm,
  belowskip=3mm,
  showstringspaces=false,
  basicstyle={\small\ttfamily},
  numbers=left,
  numbersep=1mm,
  keywordstyle=\color{blue},
  commentstyle=\color{dkgreen},
  stringstyle=\color{mauve},
  breaklines=true,
  breakatwhitespace=true,
  tabsize=3,backgroundcolor=\color{lgray},    
  columns=fullflexible,
  breakautoindent=false,
  framerule=1pt,
  columns=fixed,
  xleftmargin=0.5cm
}

\begin{document}
\justifying
\begin{center}
    % espaçamento 1.5
    \onehalfspacing
    
    \Large{INSTITUTO FEDERAL DO SUL DE MINAS GERAIS}
    
    \normalsize{Campus Poços de Caldas}
    
    \vspace{0.5cm}
    \includegraphics[scale=0.1]{img/pcs-logo.eps}
    
    \LARGE{\textbf{Método de Newton para o Aprendizado do Neurônio Artificial}}
    
    \large{\textbf{Joaquim Augusto Pomarico}$^{1}$, \textbf{Marcos Vinícius Moreira}$^{2}$, \textbf{Otávio Messias Palma}$^{3}$, \textbf{Pablo Borges dos Santos}$^{4}$, \textbf{Yuri Fernandes}$^{5}$,}
    
    \normalsize{\texttt{joaquim.cioffi@alunos.ifsuldeminas.edu.br$^{1}$}, \texttt{ marcos@alunos.ifsuldeminas.edu.br$^{2}$}, \texttt{ otavio.palma@alunos.ifsuldeminas.edu.br$^{3}$},\texttt{ pablo.borges@alunos.ifsuldeminas.edu.br$^{4}$},\texttt{ yuri@alunos.ifsuldeminas.edu.br$^{5}$}}
    
    Poços de Caldas, 09 de Dezembro de 2019
\end{center}

\noindent\textbf{Resumo:} \textit{Para uma rede neural complexa, com parâmetros muito grandes, a tarefa de encontrar o mínimo dentre vários vales contido em tal superfície é muito difícil e custosa. 
Conforme visto na disciplina, um dos métodos para encontrar tal mínimo é a utilização do Método Gradiente Descendente. Todavia, para redes neurais grandes e com mais de uma camada oculta, o método gradiente se torna menos eficiente.
Um dos problemas para uso do método gradiente em grandes redes é o fato do mesmo utilizar método de otimização de primeira ordem, assumindo sempre que a superfície da rede é um plano, sem contar com curvaturas. Uma das soluções para os pontos negativos do método gradiente é aplicação do método de Newton. 
O método de Newton é um método iterativo de aproximação, geralmente utilizado em problemas de otimização de segunda ordem e para aproximação de funções.} 

\section{Introdução} 

\begin{itemize}
    \item Utilização do Método de Newton para aprendizado do Neurônio Artifical ao invés do Método do Gradiente visto em sala.
    \item O algoritimo utilizado foi o RNA desenvolvido em sala junto do professor da disciplina, sendo que foram feitas alterações na evolução dos neurônios com a implementação do Método de Newton.
    \item Foi utilizada a lingugaem de programação Java e as IDEs VSCode e NetBeans.
    \item Breve descrição dos resultados.
\end{itemize}

\section{Fundamentação teórica}

\begin{itemize}
    \item Em análise numérica, o método de Newton, desenvolvido por Isaac Newton e Joseph Raphson, tem como objetivo estimar raízes de uma função, para isto escolhe-se uma aproximação inicial para esta e calcula-se a equação da reta tangente (através de uma derivada) da função neste ponto e a interseção dela com o eixo das abcissas, a fim de se encontrar uma melhor aproximação para a raíz. Repetindo-se o processo, cria-se um método iterativo para encontrarmos a raíz da função.
    \item Fórmula do Método de Newton:
    \begin{equation}
        x_{n+1} = x_n - \frac{f(x_n)}{f'(x_n)}, n \in \mathbb{N}
    \end{equation}
    \begin{center}
        \small{\texttt{Onde o valor inicial de x é estimado no começo da iteração.}}
    \end{center}
\end{itemize}

\section{Algoritmo}

\begin{itemize}
    \item O algoritmo é fundamentalmente o mesmo implementado em aulas com algumas pequenas alterações. 
    \item A função de ativação definida para uso foi a Limiar segundo conselho do professor. A rede neural artificial escolhida foi o Perceptron, novamente sob conselho do professor, dessa forma, não é possível lidar com mais de 2 camadas.
    \item A Rede Neural Artifical Perceptron é inicializada com o número de entradas adequado ao caso e tem a função de ativação configurada junto do número de épocas e dos vetores de aprendizado (entrada e saída esperada).
    \item Durante o aprendizado, atualiza-se os pesos da entradas a partir do valor inicial estimado, que no caso do algoritmo é gerado de forma aleatória entre 0,0 e 1,0. A cada iteração nova, calcula-se o erro até que este seja zerado, o que signifca que o valor desejado foi encontrado.
    \item Utilizando o método de Newton, diferente do Método do Gradiente, temos a possibilidade de estar dividindo por 0, dessa forma trata-se esse caso somente replicando o valor da entrada para a próxima iteração.
    \item No final da iteração mostra-se o erro total da epoca, caso seja 0 a execução foi bem sucedida, caso contrário não.
    \item No caso de uma execução bem sucedida o valor de saída indicado na linha abaixo será de acordo com os dados de entrada.
    \item A ideia de usar o algoritmo pronto veio da própria sugestão do professor que nos orientou a somente fazer as mudanças no código ao invés de refatorá-lo inteiro. Pode-se aproveitar do código pronto e focar nas alterações necessárias com mais afinco visto que não precisamos nos preocupar com a funcionalidade do código que ja fora definida como confiável previamente.
\end{itemize}

\section{Recursos de implementação}

\begin{itemize}
    \item Foi utilizado apenas a linguagem de programação Java para realização do trabalho.
    \item As IDEs NetBeans e VSCode foram utilizadas para escrever o algoritmo. O serviço GitKraken foi utilizado para versionamento do mesmo. 
    \item Nenhuma biblioteca se mostrou necessária na implementação do Algoritmo.
    \item Nenhuma base de dados se mostrou necessária na implementação do Algoritmo, apenas entradas baseadas em vetores de inteiros com suas respectivas saidas em outro vetor de inteiros.
\end{itemize}

\section{Código fonte}

\begin{itemize}
    \item Código fonte deverá ser enviado junto com este relatório.
    \item Não enviar o programa compilado.
    \item Descrever como compilar e executar o programa, se necessário.
    \item Apresentar os aspectos mais importantes da implementação. Veja o exemplo a seguir.
\end{itemize}

\lstinputlisting[language=Java]{codigo/Main.java}

\section{Conclusão}

\begin{itemize}
    \item Descrever os resultados obtidos com a aplicação do algoritmo sobre o contexto.
    \item Descrever quaisquer limitações encontradas.
    \item Descrever a contribuição do trabalho para a construção do conhecimento do grupo.
\end{itemize}

\noindent\textbf{Observações:} Serão avaliados: o relatório do trabalho prático, o código fonte e a aplicação do algoritmo sobre o contexto prático.

\begin{equation}
    NF = (NC+NR+NA)*0.1,
\end{equation}

\noindent onde:
\begin{itemize}
    \item $NF$: é a nota final do trabalho.
    \item $NR$: é a nota sobre o relatório.
    \item $NC$: é a nota da implementação, do código fonte do algoritmo e dos seus resultados sobre o contexto prático.
    \item $NA$: é a nota da apresentação e arguição oral do trabalho.
    \item $NC \in [0,10], NR \in [0,10], NA \in [0,10]$ e $NF \in [0,3]$.
\end{itemize}

\end{document}
